\documentclass[11pt]{article}
\usepackage{lmodern}
\usepackage[T1]{fontenc}
\setlength{\emergencystretch}{3em}
\usepackage[margin=1.05in]{geometry}
\usepackage{amsmath,amssymb,amsthm}
\usepackage[hidelinks]{hyperref}
\newtheorem{theorem}{Theorem}[section]
\newtheorem{corollary}[theorem]{Corollary}
\newtheorem{proposition}[theorem]{Proposition}
\newtheorem{lemma}[theorem]{Lemma}
\newcommand{\HH}{\mathbb H}
\newcommand{\Z}{\mathbb Z}
\newcommand{\Q}{\mathbb Q}
\newcommand{\R}{\mathbb R}
\DeclareMathOperator{\arsinh}{arsinh}
\DeclareMathOperator{\ord}{ord}
\title{Central derivatives of elliptic curve L-functions\\
from Fourier modes on closed geodesics}
\author{}
\date{Working draft}
\begin{document}
\maketitle

\begin{abstract}
For an elliptic curve \(E/\Q\) we construct an explicit sequence of closed
geodesic traversals on \(X_0(N)\). A single, fixed nonzero Fourier mode of the
modular differential on these traversals approximates the completed
\(L\)-function at \(1+i\tau\), with \(\tau\) inversely proportional to the
length. The absolute error is exponentially small in the length. Consequently
the decay exponent of this mode equals the analytic rank, and its leading
coefficient recovers the first nonzero central derivative with an explicit
constant. The geodesics can be taken in nested orders of a fixed real
quadratic field. The proof uses the Fricke involution, the Mellin transform
and cusp decay. Recovering the Mordell--Weil rank from the same periods
remains open.
\end{abstract}

\section{The period and the central derivative}

Let \(E/\Q\) have conductor \(N\), with normalized weight-two newform
\(f(z)=\sum_{n\ge1}a_n e^{2\pi i n z}\). Write
\[
 \omega=2\pi i f(z)\,dz,\qquad
 \Lambda_E(s)=\left(\frac{\sqrt N}{2\pi}\right)^s\Gamma(s)L(E,s),
 \qquad \Lambda_E(s)=w\Lambda_E(2-s),\quad w\in\{1,-1\}.
\]
Here \(L(E,s)\) includes the local factors at the bad primes. The Fricke
involution \(W_N z=-1/(Nz)\) satisfies \(W_N^*\omega=-w\omega\).
All heights on \(\HH\) below use the coordinate \(z=x+iy\), with hyperbolic
arclength \(|dz|/y\).

For a positive integer \(m\), put
\[
 \gamma_m=\begin{pmatrix}1+Nm^2&m\\Nm&1\end{pmatrix}\in\Gamma_0(N),
 \qquad T=T_m=\arsinh\!\left(\frac{m\sqrt N}{2}\right).
\]
The axis of \(\gamma_m\) contains \(z_*=i/\sqrt N\). Orient the axis
towards its positive real endpoint. Let \(z_m(t)\) be its unit-speed
parametrization with \(z_m(0)=z_*\), and put
\(G_m(t)=2\pi i f(z_m(t))z_m'(t)\).
The translation length of \(\gamma_m\) is \(4T_m\).
For \(h\in\Z\), define the \emph{unnormalized} Fourier integral over this
specified traversal by
\begin{equation}
 I_{m,h}=\int_{-T_m}^{3T_m} e^{i\tau_{m,h}t} G_m(t)\,dt,
 \qquad \tau_{m,h}=\frac{\pi h}{2T_m}.
 \label{eq:period}
\end{equation}
The choice of \(z_*\) fixes the phase. The Fourier coefficient with probability
arclength measure is \(I_{m,h}/(4T_m)\).

\begin{theorem}\label{thm:comparison}
For all integers \(h\) satisfying \((-1)^h=-w\), as \(m\) tends to infinity,
\begin{equation}
 I_{m,h}=-\frac{4\pi}{\sqrt N}\,
         \Lambda_E(1+i\tau_{m,h})+O_E(e^{-T_m}).
 \label{eq:comparison}
\end{equation}
The implied constant is uniform in \(h\).
\end{theorem}

\begin{corollary}\label{cor:rank}
Set \(h=1\) if \(w=1\) and \(h=2\) if \(w=-1\).
For \(r=\ord_{s=1}L(E,s)\),
\begin{equation}
 \lim_{m\to\infty} T_m^r I_{m,h}
 =-\frac{2}{r!}\left(\frac{\pi i h}{2}\right)^r L^{(r)}(E,1)\ne0 .
 \label{eq:rank}
\end{equation}
In particular, if \(w=1\) and \(L(E,1)=0\), then
\[
 \lim_{m\to\infty} T_m^2 I_{m,1}
   =\frac{\pi^2}{4}\,L''(E,1).
\]
The last identity holds also when the second derivative vanishes.
\end{corollary}

\begin{proof}[Proof of Corollary \ref{cor:rank}]
At the first nonzero derivative one has
\(\Lambda_E^{(r)}(1)=\sqrt N\,L^{(r)}(E,1)/(2\pi)\).
Taylor expansion in \eqref{eq:comparison} proves \eqref{eq:rank}. For the
final assertion, the functional equation and \(w=1\) give
\(\Lambda_E'(1)=0\). The hypothesis \(L(E,1)=0\) then also gives
\(L'(E,1)=0\).
\end{proof}

Hyperbolic Fourier expansions and their relation to geodesic periods appear
in O'Sullivan and Taylor \cite{ost}. Constantinescu and Nordentoft compare
closed-geodesic periods and vertical line integrals in their distribution
results \cite{cn}. Popa relates central Rankin \(L\)-values over real
quadratic fields to automorphic periods \cite{popa}. Our calculation concerns
one specified traversal and one low Fourier mode, with its phase retained,
as the length increases. The proof below is independent of a formula for
the square of a global toric period.

\section{The two halves of the traversal}

In the coordinate \(Z=\sqrt N z\) the fixed endpoints of \(\gamma_m\) solve
\(Z^2-m\sqrt N Z-1=0\). They are \(-e^{-T}\) and \(e^T\). The Euclidean
semicircle with these endpoints passes through \(i\). An explicit unit-speed
parametrization is
\begin{equation}
 Z(t)=\sqrt N z_m(t)=
 \frac{\lambda(u^2-1)+i u(1+\lambda^2)}{1+\lambda^2 u^2},
 \qquad \lambda=e^{-T},\quad u=e^t.
 \label{eq:param}
\end{equation}
For example, \(Z(t)\) is obtained by applying
\(v\mapsto(-e^{-T}+e^T v)/(1+v)\) to \(v=i e^{t-T}\).
The map is an orientation-preserving hyperbolic isometry, and \(Z(0)=i\).

Complex conjugation on the modular curve is induced by \(c(z)=-\bar z\),
and \(c^*\omega=\bar\omega\) since the \(a_n\) are real. The map \(W_N\)
fixes \(z_*\) and interchanges the two endpoints, so
\(W_N z_m(t)=z_m(-t)\). The antiholomorphic map
\[
 J_m(z)=m+\frac{1}{N\bar z}
\]
fixes the endpoints individually and sends \(z_*\) to \(m+z_*\) in the
positive direction on the axis. The hyperbolic distance of these points is
\(2T\), since its hyperbolic cosine is \(1+Nm^2/2=\cosh(2T)\).
Thus \(J_m z_m(t)=z_m(t+2T)\). Direct composition gives
\(J_m^2=\gamma_m\), proving the asserted translation length.
The traversal is primitive even in \(\mathrm{PSL}_2(\Z)\). Indeed a
positive translation along the same oriented axis has an integral determinant-one
representative of trace \(t>2\) and upper right entry \(b>0\). Its fixed-point
equation must be proportional to \(Nz^2-Nmz-1=0\), hence its lower left entry
is \(Nb\) and its diagonal difference is \(Nmb\). Thus
\(t^2=(Nmb)^2+4Nb^2+4\), with \(t\equiv Nmb\pmod 2\). If \(0<b<m\),
then \(Nmb<t<Nmb+2\), an impossibility. The trace increases with \(b>0\),
and \(\gamma_m\) realizes \(b=m\).

Translation by \(m\) preserves \(\omega\), and \(J_m\) is this translation
composed with \(W_N c\). Pullback of the differential therefore gives
\begin{equation}
 G_m(-t)=wG_m(t),\qquad
 G_m(t+2T)=-w\overline{G_m(t)}.
 \label{eq:symmetries}
\end{equation}
In the first identity the sign from reversing the parameter is essential.
On splitting \eqref{eq:period} at \(T\) and using
\(e^{2iT\tau_{m,h}}=(-1)^h=-w\), we obtain the exact identity
\begin{equation}
 I_{m,h}=2\int_{-T}^T e^{i\tau_{m,h}t}\operatorname{Re}G_m(t)\,dt .
 \label{eq:half}
\end{equation}

\section{Comparison with the vertical Mellin integral}

On the vertical geodesic \(z=i e^t/\sqrt N\), the corresponding coefficient
of the differential is real and equals
\[
 G_\infty(t)=-\frac{2\pi}{\sqrt N}\,e^t f(i e^t/\sqrt N).
\]
It satisfies \(G_\infty(-t)=wG_\infty(t)\) by the same Fricke calculation.

\begin{lemma}\label{lem:error}
As \(m\to\infty\),
\[
 \int_{-T}^{T}|G_m(t)-G_\infty(t)|\,dt+
 \int_{|t|>T}|G_\infty(t)|\,dt=O_E(e^{-T}).
\]
\end{lemma}

\begin{proof}
Both absolute integrands are even by \eqref{eq:symmetries} and the vertical
symmetry. For \(0\le t\le T\), put \(u=e^t\), so that \(1\le u\le1/\lambda\).
Equation \eqref{eq:param} gives \(\operatorname{Im}Z(t)\ge u/2\). The same
formula and its derivative give the uniform estimates
\[
 |Z(t)-iu|+|Z'(t)-iu|\ll \lambda u^2,\qquad |Z'(t)|\ll u.
\]
Indeed the real numerator of \(Z'(t)\) over
\((1+\lambda^2u^2)^2\) is \(2\lambda u^2(1+\lambda^2)\), and its
imaginary numerator is
\(u(1+\lambda^2)(1-\lambda^2u^2)\).
On the segment joining \(Z(t)/\sqrt N\) to \(iu/\sqrt N\), the Fourier
expansion at infinity bounds \(|f|+|f'|\) by \(C_Ee^{-c_Eu}\), with
\(C_E,c_E>0\). It follows that
\[
 |G_m(t)-G_\infty(t)|\ll_E
 \lambda(u^2+u^3)e^{-c_Eu}.
\]
The integral against \(du/u\) is \(O_E(\lambda)\). The same Fourier bound
on the vertical ray gives
\(\int_T^\infty|G_\infty(t)|\,dt\ll_E e^{-c_E e^T}\), after adjusting
\(c_E\) if needed.
\end{proof}

\begin{proof}[Proof of Theorem \ref{thm:comparison}]
Since \(|e^{i\tau t}|=1\) for real \(\tau\), \eqref{eq:half} and
Lemma \ref{lem:error} replace \(I_{m,h}\) by
\(2\int_\R e^{i\tau_{m,h}t}G_\infty(t)\,dt\) with error \(O_E(e^{-T})\).
The substitution \(y=e^t/\sqrt N\) and the Mellin transform of \(f(iy)\)
give
\[
 \int_\R e^{i\tau t}G_\infty(t)\,dt
  =-2\pi N^{i\tau/2}\int_0^\infty y^{i\tau}f(iy)\,dy
  =-\frac{2\pi}{\sqrt N}\Lambda_E(1+i\tau).
\]
The Mellin identity follows first in the half-plane of absolute convergence
of the Dirichlet series by integrating the Fourier expansion. Cusp decay at
both ends continues the integral to every complex Mellin parameter.
\end{proof}

\section{Nested orders in a fixed real quadratic field}

The endpoint in the original \(z\)-coordinate satisfies
\(Nz^2-Nmz-1=0\). Its primitive quadratic discriminant is
\begin{equation}
 \Delta_m=N^2m^2+4N=4N\cosh^2 T_m.
 \label{eq:disc}
\end{equation}
This number is positive and nonsquare since \(m^2\Delta_m=(2+Nm^2)^2-4\) lies strictly
between two consecutive squares. For either endpoint \(\alpha\), the
multiplier order of \(\Z+\Z\alpha\) is \(\Z+\Z N\alpha\).
Indeed any multiplier is \(a+b\alpha\) with \(a,b\in\Z\). Its value on
\(\alpha\) belongs to the lattice exactly when \(N\mid b\), by
\(\alpha^2=m\alpha+1/N\). This order has discriminant \(\Delta_m\).

\begin{proposition}\label{prop:nesting}
Fix \(m_*\ge1\) and put \(T_*=T_{m_*}\). For each odd positive integer
\(M\), set \(m(M)=2\sinh(MT_*)/\sqrt N\). Then \(m(M)\) is a positive
integer and \(T_{m(M)}=MT_*\). All the multiplier orders belong to
\(K=\Q(\sqrt{\Delta_{m_*}})\). If \(b(M)\) is their conductor in the
maximal order of \(K\), then
\begin{equation}
 \frac{b(M)}{b(1)}=\frac{\cosh(MT_*)}{\cosh T_*}.
 \label{eq:conductor}
\end{equation}
When \(M\mid M'\) and both are odd, \(b(M)\mid b(M')\). In particular,
for any odd prime \(q\) the traversals indexed by \(M=q^j\) have lengths
\(4q^j T_*\), and their multiplier orders form a descending chain.
\end{proposition}

\begin{proof}
For \(t=T_*\), let \(x=4\sinh^2t=Nm_*^2\). As functions of positive odd
\(M\), the sequences \(S_M=\sinh(Mt)/\sinh t\) and
\(C_M=\cosh(Mt)/\cosh t\) both satisfy
\(A_{M+4}=(x+2)A_{M+2}-A_M\). Their first two terms are
\((1,x+3)\) and \((1,x+1)\), respectively. They are therefore positive
integers. In particular \(m(M)=m_* S_M\). Equation \eqref{eq:disc} now
gives \(\Delta_{m(M)}=\Delta_{m_*}C_M^2\), proving the common field and,
using the discriminant-conductor formula for quadratic orders,
\eqref{eq:conductor}. For \(M'=ML\), applying the same recurrence at
\(t'=MT_*\) gives \(\cosh(Lt')/\cosh t'\in\Z\), since
\(4\sinh^2 t'=Nm(M)^2\). This proves divisibility of the conductors.
\end{proof}

\section{The arithmetic question}

The unweighted integral over each of these traversals is zero. In fact
\[
 \gamma_m=
 \begin{pmatrix}1&m\\0&1\end{pmatrix}
 \begin{pmatrix}1&0\\Nm&1\end{pmatrix}.
\]
The two factors are parabolic elements fixing cusps. Integration of a
holomorphic differential defines an additive period homomorphism on
\(\Gamma_0(N)\). It vanishes on parabolic elements because a loop around
a cusp bounds a disk on \(X_0(N)\). The nonzero modes in
Corollary \ref{cor:rank} retain information lost by that unweighted integral.

For a fixed \(E\), equation \eqref{eq:rank} gives
\(\lim_{m\to\infty}\log |I_{m,h}|/\log T_m=-r\).
The rank formulation of Birch and Swinnerton-Dyer asks whether \(r\)
equals \(\operatorname{rank}E(\Q)\) \cite{wiles}. An independent
arithmetic determination of this decay exponent in terms of rational
points would therefore suffice. The comparison proved here determines
the exponent analytically. Even in the nested-order family the
character \(e^{i\tau_{m,h}t}\) varies on the real geodesic, while the
multiplier orders have the conductors in \eqref{eq:conductor}.
Nesting by itself gives no relation between the period and a lattice
of rational points.

\begin{thebibliography}{9}
\bibitem{cn} P. Constantinescu and A. C. Nordentoft,
\emph{Non-vanishing of geodesic periods of automorphic forms},
\href{https://doi.org/10.1007/s00039-025-00715-z}{Geom. Funct. Anal. (2025)}.
\bibitem{ost} C. O'Sullivan and K. Taylor,
\emph{Hyperbolic Fourier coefficients of Poincar\'e series},
\href{https://doi.org/10.1007/s11139-016-9772-6}{Ramanujan J. (2016)}.
\bibitem{popa} A. A. Popa,
\emph{Central values of Rankin \(L\)-series over real quadratic fields},
\href{https://doi.org/10.1112/S0010437X06002259}{Compos. Math. (2006)}.
\bibitem{wiles} A. Wiles, \emph{The Birch and Swinnerton-Dyer conjecture},
\href{https://www.claymath.org/wp-content/uploads/2022/05/birchswin.pdf}
{Clay Mathematics Institute, official problem description}.
\end{thebibliography}
\end{document}
